% Secao 5: Mao na Massa (Laboratorios de Ponta, Racks e Projetos)

% Slide 12: Laboratorios de Ponta
\begin{frame}{Mão na Massa: Laboratórios de Padrão Industrial}{\faMicrochip\ Nada de ficar só olhando pro quadro: você aprende fazendo com equipamentos reais}
  \begin{columns}[t]
    \begin{column}{0.48\textwidth}
      \begin{greencard}{\faNetworkWired\ Lab de Redes \& Telecomunicações}
        \tiny
        \begin{itemize}\setlength{\itemsep}{0.18em}
          \item \textbf{Racks Industriais Reais:} Os mesmos gabinetes e estruturas usados por grandes provedores e data centers.
          \item \textbf{Equipamentos de Ponta:} Switches gerenciáveis, roteadores Gigabit Dual-Band e centrais VoIP.
          \item \textbf{Bancada de Ferramental:} Alicates de crimpagem, testadores de cabos de rede e fusão de fibra óptica.
          \item \textbf{Simulação \& Prática:} Aprenda a capturar pacotes e auditar o tráfego com Wireshark e Cisco Packet Tracer.
        \end{itemize}
      \end{greencard}
    \end{column}
    
    \begin{column}{0.48\textwidth}
      \begin{bluecard}{\faDesktop\ Laboratórios de Informática 1 e 2}
        \tiny
        \begin{itemize}\setlength{\itemsep}{0.18em}
          \item \textbf{Mais de 60 Computadores:} Máquinas modernas e rápidas para cada aluno praticar de forma individual.
          \item \textbf{Conforto Térmico Total:} Ambientes 100\% climatizados com ar-condicionado e estações ergonômicas.
          \item \textbf{Internet Ultrarrápida Federal:} Conexão de fibra em alta velocidade para downloads e pesquisas.
          \item \textbf{Acessibilidade Completa:} Postos adaptados para Pessoas com Deficiência (PcD).
        \end{itemize}
      \end{bluecard}
    \end{column}
  \end{columns}
  
  \vspace{0.08cm}
  \begin{statcard}
    \tiny \textcolor{IFCEDarkGreen}{\faCheckCircle\ \textbf{Diferencial Único:} Na escola comum o computador só serve para digitar trabalhos; no IFCE você aprende como ele funciona por dentro e como fazê-lo conversar com o mundo!}
  \end{statcard}
\end{frame}

% Slide 13: Projetos Reais, Robotica e TECTEL
\begin{frame}{Projetos Reais, Robótica e Eventos Científicos}{\faFlask\ Desenvolva ideias inovadoras e apresente soluções para a sua comunidade}
  \begin{columns}[t]
    \begin{column}{0.48\textwidth}
      \begin{purplecard}{\faRobot\ Robótica \& Internet das Coisas (IoT)}
        \tiny
        \begin{itemize}\setlength{\itemsep}{0.18em}
          \item \textbf{Dispositivos Inteligentes:} Conecte placas Arduino e ESP32 com sensores à rede Wi-Fi.
          \item \textbf{Automação Residencial e Rural:} Sistemas de irrigação automática, sensores de temperatura e controle remoto por celular.
          \item \textbf{Projetos Integradores:} Equipes de alunos criando protótipos funcionais para desafios reais de Tauá.
        \end{itemize}
      \end{purplecard}
    \end{column}
    
    \begin{column}{0.48\textwidth}
      \begin{greencard}{\faCalendar\ Eventos \& Feiras de Ciência}
        \tiny
        \begin{itemize}\setlength{\itemsep}{0.18em}
          \item \textbf{TECTEL (Semana da Tecnologia):} O maior evento tech da região, com oficinas práticas, palestras e minicursos.
          \item \textbf{Hackathons \& Maratonas:} Desafios estimulantes para criar soluções de programação em equipe.
          \item \textbf{Viagens Técnicas:} Visitas guiadas a provedores de internet, emissoras e empresas de tecnologia.
        \end{itemize}
      \end{greencard}
    \end{column}
  \end{columns}
  
  \vspace{0.08cm}
  \begin{center}
    \scriptsize \textcolor{IFCEDarkGreen}{\textbf{\faTrophy\ Destaque Regional: Os estudantes do IFCE Tauá frequentemente conquistam medalhas em feiras científicas estaduais e nacionais!}}
  \end{center}
\end{frame}
